\section{Data and Preprocessing}
\label{sec:data}

\subsection{AI-READI multimodal panel}

The AI-READI dataset provides a 5-minute multimodal panel assembled from CGM (Dexcom),
wrist-worn wearables (HR, steps, sleep stages, skin-conductance stress), static clinical
features (HbA1c, BMI, medication flags, demographics, clinical site), and a predicted meal
flag derived from a trained meal-detection model.
The panel is stored in a single enriched Parquet file
(\texttt{final\_multimodal\_dataset\_*}), with a companion
\texttt{participant\_static\_features.parquet} providing per-participant covariates.

\subsection{Cohort selection}

A four-stage cohort selection procedure selects participants with sufficient CGM continuity
for streaming training:
\begin{enumerate}
  \item \textbf{CGM coverage}: participants with $\ge 49$~h of continuous CGM segments
        (minimum \texttt{clean\_min\_segment\_hours}).
  \item \textbf{Stratified split}: 70/15/15 participant-level holdout stratified on
        \texttt{participants\_study\_group} (healthy, pre-diabetes lifestyle-controlled,
        oral/injectable medication, insulin-dependent), seed 42.
  \item \textbf{Segment boundary enforcement}: no window or streaming state crosses a
        (\texttt{participant\_id}, \texttt{segment\_id}) boundary.
        The streaming unit is \texttt{participant\_id\_plus\_segment\_id}.
  \item \textbf{Scaler fit on train only}: all normalization and per-horizon robust scales are
        fit on train participants and applied to validation/test without leakage.
\end{enumerate}

Table~\ref{tab:cohort} summarizes the cohort characteristics.
Full participant counts per split are \TODO{insert from split CSV after full-cohort run}.

\begin{table}[h]
\centering
\caption{AI-READI cohort characteristics used in this study.  Values are as
  reported in the published dataset description.}
\label{tab:cohort}
\IfFileExists{tables/generated/cohort_characteristics.tex}{%
  \begin{tabular}{lc}
\toprule
Characteristic & Value \\
\midrule
Split & Participant-held-out 70/15/15 \\
Sampling interval & 5 min \\
Forecast horizon & 12 steps / 60 min \\
Train streams & 1,849 \\
Validation streams & 392 \\
Test streams & 381 \\
Test forecast anchors & 155,113 \\
Timed insulin dose / IOB & Not available in AI-READI \\
\bottomrule
\end{tabular}
%
}{%
  \begin{tabular}{lc}
  \toprule
  Characteristic & Value \\
  \midrule
  Study groups & healthy, pre-diabetes, oral/injectable med, insulin-dependent \\
  CGM device & Dexcom \\
  Sampling interval & 5 min \\
  Wearables & HR, steps, sleep, stress \\
  Static features & HbA1c, BMI, medications, demographics, site \\
  Meal proxy & predicted meal flag (meal-detection model) \\
  Insulin covariate & \texttt{med\_insulin} static flag only (no timed dose/IOB) \\
  \bottomrule
  \end{tabular}%
}
\end{table}

\subsection{Predicted meal flag}

The \texttt{predmeal\_flag} is produced by an upstream meal-detection model applied to the
CGM slope and contextual covariates.
It is a \emph{noisy proxy}: it may fire retrospectively (flagging a glucose rise that already
occurred), miss low-carbohydrate meals, or fire spuriously during hyperglycemic episodes.
It is used in this work only as a \texttt{meal\_proxy} scenario covariate
(\S\ref{sec:proxy}), not as a ground-truth meal time or carbohydrate quantity.
Any scenario evaluated with \texttt{predmeal\_flag = 1} should be interpreted as a
retrospective simulation, not a prospective intervention.

\subsection{Gap handling and segment resets}

CGM gaps longer than 30~minutes trigger a new segment.  The streaming state is reset to the
static-initialized $h_0$ at each segment boundary.  Windows and training chunks do not cross
segment boundaries.  The residual-current target transform anchors every forecast to the
last observed glucose, ensuring that a fresh segment cold-starts at full accuracy without
requiring warm-up steps.
